\section{AI Bridge Approach}

AI Bridge is used as a local execution and orchestration layer. The intended
workflow is shown in Figure~\ref{fig:pipeline}.

\begin{figure}[ht]
\centering
\begin{tikzpicture}[
  node distance=1.4cm,
  block/.style={draw, rounded corners, align=center, minimum width=2.8cm, minimum height=0.8cm},
  arrow/.style={->, thick}
]
\node[block] (inputs) {Environment\\Evidence};
\node[block, right=of inputs] (watcher) {AI Bridge\\Watcher};
\node[block, right=of watcher] (agents) {Specialized\\AI Agents};
\node[block, below=of agents] (scoring) {CVSS Env.\\Scoring};
\node[block, left=of scoring] (audit) {Evidence\\Audit Trail};
\node[block, left=of audit] (report) {Report +\\Review};

\draw[arrow] (inputs) -- (watcher);
\draw[arrow] (watcher) -- (agents);
\draw[arrow] (agents) -- (scoring);
\draw[arrow] (scoring) -- (audit);
\draw[arrow] (audit) -- (report);
\draw[arrow] (watcher) |- (audit);
\end{tikzpicture}
\caption{Evidence-based CVSS Environmental assessment workflow using AI Bridge.}
\label{fig:pipeline}
\end{figure}

We define four agent roles:
\begin{enumerate}[noitemsep]
  \item Network/Topology Analyst;
  \item PCI-DSS Scope Reviewer;
  \item CVSS Environmental Analyst;
  \item Evidence Consolidator.
\end{enumerate}

The agents do not simply output a score. Each metric value must reference
evidence, such as asset inventory, business impact classification, firewall
rules, topology, or vulnerability metadata.
